\uuid{qemP}
\titre{Calcul de dérivée}
\auteur{bodin ; hanani ; mzari}
\organisation{exo7}
\licence{CC BY-NC-SA 4.0}
\datecreate{2020-01-20}
\niveau{L1}
\module{Analyse}
\chapitre{Dérivabilité des fonctions réelles}
\sousChapitre{Autre}
\difficulte{1}
\theme{Dérivées, calcul de dérivée}
\exo7qcmid{lille-1.derivees.005}

\contenu{
\begin{question}
\qclassification{124.00}
\qtags{motcle=calcul de dérivée}
Quelles sont les bonnes réponses ?
\begin{answers}
    \good{La dérivée de $f(x)=(2x+1)^2$ est $f'(x)=4(2x+1)$.}
    \bad{La dérivée de $f(x)=(2x+1)^2$ est $f'(x)=2(2x+1)$.}
    \bad{La dérivée de $f(x)=\mathrm{e}^{x^2-2x}$ est $f'(x)=2\mathrm{e}^{x^2-2x}$.}
    \good{La dérivée de $f(x)=\mathrm{e}^{x^2-2x}$ est $f'(x)=2(x-1)\mathrm{e}^{x^2-2x}$.}
\end{answers}
\begin{explanations}
De manière plus générale, $(u^n)'=nu^{n-1}u'$ et $(\mathrm{e}^v)'=v'\mathrm{e}^v$. Il suffit de prendre $u=2x+1$, $n=2$ et $v=x^2-2x$.
\end{explanations}
\end{question}
}
